\section{Relative laws, normal form, and the retained smooth invariant}
\label{sec:relative}
The asymmetric inverse above is a short-record theorem. It neither
requires nor replaces the function-valued long-record result. We state
the distinction precisely and include the analytic comparison in a
form with an explicit physical chart convention.

\subsection{The smooth physical law}
For $j$ alternating flights starting at contact $0$, write
$E_j=W_j-jg$, $p=j\bmod2$, and
$b_j=-W_{j,uv}/d_j^0$. Define
\[
 a_b=\frac{\sqrt{c_0c_1}\sinh\gamma}{gc_{1-b}},\qquad
 q_p=\sqrt{a_0a_p},\qquad d_j^0=\frac{q_p}{\sinh(j\gamma)}.
\]
The complete smooth result proved in \cite{Companion}, in the sections
on the half-line boundary layers and differentiated operators, is as
follows. On fixed common neighborhoods there are positive smooth $B_b$
and strictly convex smooth $S_b$ with
\[
 B_b(0)=1,\quad S_b(0)=S'_b(0)=0,\quad S_b''(0)=a_b,
\]
such that for each fixed derivative order $k$,
\begin{equation}\label{eq:smoothrelative}
 \|E_j-S_0\oplus S_p\|_{C^k}
       +\|b_j-B_0\otimes B_p\|_{C^k}\le C_k\tau^j,
                    \qquad 0<\tau<1.
\end{equation}
The norms include fixed mixed geometric derivatives on the compact
local families specified there. Normalization of the exact cofactor
identity precedes trace-norm comparison of the two end perturbations;
an uncontrolled absolute error is never divided by $d_j^0$.
The common Morse domain then gives, down to the onset,
\begin{equation}\label{eq:integratedrelative}
 \left\|\frac{2A\sinh(j\gamma)}{d^2}P_{j,0}(d)
       -\mathcal F_{0p}(d)\right\|_{C^k([0,d_*])}
                              \le C_k\tau^j,
\end{equation}
where
\begin{equation}\label{eq:limitlaw}
 \mathcal F_{bc}(d)=\frac{\sqrt{a_ba_c}}{\pi d^2}
       \int(d-S_b(u)-S_c(v))_+B_b(u)B_c(v)\dd u\dd v.
\end{equation}
The corresponding result holds with the other starting contact.
These are smooth, parameter-differentiated statements on a common
physical collar, not merely Taylor-series identities.

The full proofs and all earlier applications remain in the companion
volume, including the first-hit localization, common-domain integration,
operator estimates, and both parities. The present theorem does not
use an asymptotic approximation as a substitute for the exact
finite-flight density: Lemma~\ref{lem:flux} supplied that density
directly.

\subsection{The analytic comparator in physical coordinates}
Suppose analytic canonical coordinates put a local return map in the
form
\begin{equation}\label{eq:normalform}
 N(s,p)=(\Delta(sp)s,\Delta(sp)^{-1}p),\quad
                   \Delta(0)=\lambda\in(0,1).
\end{equation}
This is the local hyperbolic normal-form setting recalled in
\cite[Section 2]{DSKL}; its local use is distinct from the hypotheses
of that paper's global inverse theorem. In this subsection the charts
are supplied, not constructed for an arbitrary smooth family.

\begin{proposition}[Relative projection formula]\label{prop:normalform}
For the mixed endpoints $s_0=s$, $p_n=t$, the small solution satisfies
\begin{equation}\label{eq:mixed}
 I=st\Delta(I)^n,\quad p_0=t\Delta(I)^n,\quad
 s_n=s\Delta(I)^n,\quad
 T_n:=\partial_t p_0=
 \frac{\Delta(I)^n}{1-nI\Delta'(I)/\Delta(I)}.
\end{equation}
On a fixed smaller neighborhood, for every fixed $k$ and some
$\rho\in(\lambda,1)$,
$\|\lambda^{-n}T_n-1\|_{C^k}\le C_k\rho^n$.
Let the physical canonical charts be
\[
 (u,P)=(U(s,p),\mathcal R(s,p)),\quad
 (v,Q)=(V(q,t),\mathcal Z(q,t)),\quad U_s(0)V_t(0)\ne0,
\]
with generating convention $P=-W_{n,u}$, $Q=W_{n,v}$.
For $\Phi_n(s,t)=(U(s,p_0),V(s_n,t))$ and
$\mathcal D_n=\det D\Phi_n$, the exact physical twist is
\begin{equation}\label{eq:projection}
 |W_{n,uv}|=\left|\frac{T_n}{\mathcal D_n}\right|
                       \circ\Phi_n^{-1},\qquad
 \mathcal D_n(0)=U_s(0)V_t(0)-\lambda^{2n}U_p(0)V_q(0).
\end{equation}
The inverses exist on one physical box for all sufficiently large $n$.
If $\sigma(u)$ and $\theta(v)$ invert $U(s,0)$ and $V(0,t)$,
the normalized physical twist converges there in every fixed $C^k$
norm to
\begin{equation}\label{eq:projectionfactors}
 \frac{U_s(0)}{U_s(\sigma(u),0)}
             \frac{V_t(0)}{V_t(0,\theta(v))}.
\end{equation}
\end{proposition}
\begin{proof}
Choose $\lambda<\mu<\rho<1$ and a complex disk where $\Delta$
is holomorphic, nonzero, and bounded by $\mu$. On a sufficiently small
endpoint bidisk the map $I\mapsto st\Delta(I)^n$ is a contraction
uniformly in $n$, since $\sup_n n\mu^{n-1}<\infty$.
It yields a holomorphic solution $I=O(\mu^n)$. Invariance of $sp$
gives the middle two formulas in \eqref{eq:mixed}; the intermediate
points $(s\Delta(I)^i,t\Delta(I)^{n-i})$ stay in the chart.
Implicit differentiation gives its last formula, valid also at $s=0$.
The analytic logarithm gives
\[
 \Delta(I)^n/\lambda^n=1+O(n\mu^n),\qquad
 1-nI\Delta'(I)/\Delta(I)=1+O(n\mu^n).
\]
These are estimates of normalized quantities. Cauchy bounds on a
smaller fixed bidisk and $n\mu^n\le C\rho^n$ give every fixed
derivative estimate.

It follows that $\Phi_n$ tends in each fixed derivative norm to
$(U(s,0),V(0,t))$. Transversality, a smaller physical box, and the
inverse-function contraction give inverse maps on that same box,
converging with their derivatives. Pulling back
$\dd u\wedge\dd P=\dd s\wedge\dd p_0$ to $(s,t)$ gives
$(-W_{n,uv})\mathcal D_n=T_n$. Its absolute value is
\eqref{eq:projection}; the derivative at zero gives the determinant
there. Dividing by the exact origin value and using convergence of the
inverse maps gives \eqref{eq:projectionfactors}. Each scalar derivative
has constant sign on the chosen box, so its normalized ratio is positive.
\end{proof}

The same-section specialization needs an identification of charts, not
just an equality of section names. Choose the same canonical chart at
both ends with $(q,t)$ in the same stable--unstable order as $(s,p)$,
and the same physical transverse and momentum conventions. Then
$V(q,t)=U(q,t)$, $V_t(0)=U_p(0)$, and $V_q(0)=U_s(0)$.
In particular $U_s(0)U_p(0)\ne0$, and the invariant formula becomes
\begin{equation}\label{eq:samesection}
 |W_{n,uv}(0)|=
 \frac{\lambda^n}{|U_s(0)U_p(0)|(1-\lambda^{2n})}.
\end{equation}
A reversal or interchange of a chart must instead be substituted into
\eqref{eq:projection}. For the two-step billiard return, $j=2n$ and
$\lambda=e^{-2\gamma}$. This recovers the exact hyperbolic-sine
normalization rather than only its leading exponential.

Applying the analytic comparison and \eqref{eq:smoothrelative} to the
same analytic physical problem identifies their normalized endpoint
factors by uniqueness of the limit and evaluation on the two axes.
Thus the analytic product phenomenon is not itself the new inverse
mechanism. The odd blocks in \eqref{eq:oddblock} use physical moment
information that the normal form by itself does not supply.

\subsection{The information retained after transverse symmetrization}
Define a locally finite measure on $[0,d_*]$ by
\[
 \nu_b=(S_b)_*\{\sqrt{a_b}B_b(u)\dd u\}.
\]
Near zero its density has the form
$\sqrt2\,x^{-1/2}\mathcal V_b(x)$, where
$\mathcal V_b$ is smooth and $\mathcal V_b(0)=1$.
The two inverse branches of the strictly convex $S_b$ are added in
this density. If $H_{bc}(d)=d^2\mathcal F_{bc}(d)$, then
\[
 H_{bc}(d)=\frac1\pi\int(d-x-y)_+\nu_b(\dd x)\nu_c(\dd y).
\]
Hence the smooth count law is a symmetrized energy pushforward.
The Volterra uniqueness and compatibility proof in \cite{Companion}
recovers the entire $\mathcal V_0,\mathcal V_1$ from the two
diagonal laws and predicts the mixed law. It is not an identification
of two smooth functions from their Taylor series.

The same volume gives the weighted Abel inverse norm, the two-derivative
gain from integrated flux, and the separately charged smooth-class
calibration procedure. Its full-profile preparation bound is sufficient,
not a matching minimax theorem. The limiting even-contact inverse and
its finite-jet equivalence to probability-only two-flight coefficients
remain valid with their original hypotheses. They are distinct from
the all-degree signed inverse proved here. In particular a classical
analytic normal form, a smooth symmetrized profile, and an asymmetric
contact graph are three different mathematical objects.
